\section{Background \& Related Work}
Naturally, the pursuit of structured generated outputs has two distinct branches
or "paradigms" to approach the problem: inference-time constraint decoding or
training-time weight alignment. Thus, we must analyze the computational
trade-offs in these approaches, and identify the critical latency-verification
gaps that necessitate a dual-objective framework like this.

% Review FSM based masking (outlines, JSON schemas, etc...)
  \subsection{Inference-Time Constraint Mechanisms}
  As aforementioned, the standard method for enforcing structural validity relies
  on projecting the generative vocabulary onto a formal grammar during the
  decoding step in the transformer. Frameworks such as Outlines~\cite{willard2023efficient} and
  LMQL~\cite{beurer2023prompting} utilize FSMs to compute a boolean mask over
  the vocabulary space $V$ at each timestep $t$. By mapping invalid token
  probabilities to $-\infty$, these systems can guarantee adherence to
  CFGs, such as standard JSON syntax, with negligible
  latency overhead.

  However, the memoryless nature of FSM-based decoding renders it fundamentally
  incapable of enforcing Context-Sensitive Invariants (CSIs). Complex
  domain-specific languages (DSLs) and API schemas frequently demand cross-field
  arithmetic bounds, dynamic type resolution, or referential integrity (e.g.,
  referencing a dynamically instantiated variable in a subsequent payload block).
  These constraints must therefore be validated \textit{post-hoc}, requiring
  regeneration when violations occur and increasing system latency. Attempting to
  encode such semantics directly within an FSM would result in a combinatorial
  state-space explosion, rendering the approach computationally infeasible for
  richly typed DSLs.

  Consequently, any viable formal verification framework must decouple the
  enforcement of CFG structure from higher-order semantic constraints.

  \subsection{Semantic-Aware Decoding}
  To enforce operational semantics, recent literature has explored integrating
  symbolic validation directly into the decoding process.

  PICARD~\cite{scholak2021picard}, for instance, parses the top-$k$ beams into
  Abstract Syntax Trees (ASTs) during decoding and prunes candidates that violate
  database schemas. While this elevates constraint enforcement from syntax to
  semantics, the neural model itself remains unaware of these constraints and
  therefore frequently proposes invalid distributions, forcing the system into
  repeated rejection cycles.

  Similarly, Syncromesh~\cite{poesia2022synchromesh} employs completion engines that query an external
  abstract machine to compute admissible continuations during generation. While
  mathematically rigorous, the continuous token-level interaction between the
  neural decoder and external symbolic state tracker introduces significant overhead.

  To mitigate these costs, Kogler et al.~\cite{kogler2026towards} decouple semantic verification
  from token-level decoding by delegating constraint evaluation to discrete
  structural boundaries. When a completed element violates a constraint, the
  system performs a rollback and regenerates the element.
  This strategy preserves the rigor of symbolic validation while minimizing
  synchronization overhead.

  % How does LoRA work and is it really efficient?
  \subsection{Parameter-Efficient Fine-Tuning (PEFT)}
  Parallel to decoding research, efforts have also focused on shifting the
  model's predictive density ($P_\theta(Y|X)$) towards a target DSL through
  training-time alignment.

  Gorilla~\cite{patil2024gorilla}, an LLM connected with massive APIs, demonstrated that Supervised Fine-Tuning (SFT)
  over API call traces can drastically reduce structural hallucination rates.
  Furthermore, Low-Rank Adaptation (LoRA)~\cite{hu2022lora} has proven highly effective for
  this task. By freezing pre-trained weights and injecting trainable low-rank
  update matrices, LoRA enables models to internalize the structural patterns of
  specific DSLs with minimal computational overhead.

  However, purely neural approaches cannot provide formal guarantees. Even if
  fine-tuning pushes a model to extremely high accuracy, the inherent
  probabilistic nature of transformer decoding implies that invalid outputs may
  still occasionally occur. Consequently, PEFT methods are best viewed as an
  optimization mechanism rather than a standalone solution for enforcing strict
  operational semantics.

  % Approaches combining both?
  \subsection{Dual-Objective Architectures}
  Finally, several recent systems attempt to unify neural predictive models with symbolic
  reasoning engines, combining training-time alignment with inference-time
  verification.

  AlphaGeometry~\cite{trinh2024solving} couples a neural language model trained on synthetic
  reasoning traces with a deterministic symbolic deduction engine. In this
  architecture, the neural model acts as a heuristic prior that navigates the
  combinatorial search space, while the symbolic system provides the ground truth
  of semantic correctness.

  However, these systems typically structure the interface between neural and
  symbolic components as an iterative propose-and-verify loop~\cite{jiang2022draft}. The language model
  generates a complete candidate sequence, which is then passed to a symbolic
  engine for evaluation. If a violation is detected, the engine returns an error
  trace and the model must regenerate the output~\cite{jiang2022draft}.

  This macro-level architecture introduces three key limitations:
  \begin{itemize}
  \item Latency: repeated hand-offs between the GPU-bound neural model and the
  CPU-bound symbolic verifier significantly degrade inference throughput.
  \item Redundant compute: the model frequently generates large invalid subtrees
  before violations are detected.
  \item Lack of quantification: these systems cannot formally measure how neural
  policy alignment reduces the symbolic engine's workload.
  \end{itemize}

  A closely related approach, NeSyS~\cite{zhao2026neuro}, alternates training between a
  symbolic world model and a neural generator, allowing symbolic rules to
  constrain the model's output probability distribution. Similarly, Kogler
  et al.~\cite{kogler2026towards} proposed a neuro-symbolic DSL generation framework that combines
  token-level syntactic validation with an element-level logical reasoning
  constraint solver.

  While these systems demonstrate that symbolic engines can constrain neural
  outputs, they still treat the neural model and constraint solver as largely
  static partners. Their architectures optimize the inference engine to handle
  constraint solving, but accept the computational overhead of verification as a
  fixed cost. As a result, they do not leverage the neural policy itself to
  proactively reduce the workload of the symbolic verifier.